% Part: first-order-logic
% Chapter: syntax
% Section: intro-syntax

\documentclass[../../../include/open-logic-section]{subfiles}

\begin{document}

\olfileid{fol}{syn}{itx}

\olsection{تعارف}

پہلے درجے دی منطق دا نظریہ تے
ماورائے نظریہ ودھاؤن لئی سانوں
پہلاں اوہدی لکھتاں دی نحو تے
معنویات تعریف کرنیاں پین
گیاں۔ پہلے درجے دی منطق
دیاں لکھتاں ٹرم تے
!!{formula}s ہوندیاں نیں۔
ٹرم !!{variable}s،
!!{constant}s تے
!!{function}s توں
بندے نیں۔ دوجے پاسے،
!!^{formula}s،
!!{predicate}s نوں
ٹرماں نال ملا کے
بندے نیں (ایہ سب
توں چھوٹے، ”ایٹمی“
!!{formula}s ہوندے
نیں)، تے فیر ایٹمی
!!{formula}s توں
منطقی رابطے تے
مسوّر ورت کے ہور
گنجل دار فارمولے
بنا سکدے آں۔
تشکیل دے قاعدے
بیان کرن دے کئی
طریقے نیں؛ اسیں
صرف اک ممکن
طریقہ دیندے آں۔
ہور نظام وکھریاں
علامتاں چُنن گے،
وکھرے منطقی
رابطیاں نوں
ابتدائی منّن
گے، تے قوسین
وکھرے ڈھنگ
نال ورتن گے
(یا بالکل نہیں
ورتن گے، جیویں
پولش علامت
نگاری وچ)۔
پر سارے طریقیاں
وچ سانجھی گل
ایہ اے کہ تشکیل
دے قاعدے ٹرماں
تے !!{formula}s
دے سیٹ دی تعریف
\emph{استقرائی طور تے}
کردے نیں۔ جے
ایہ کم ٹھیک کیتا
جائے، تاں ہر لکھت
تشکیل دے قاعدیاں
مطابق بنیادی طور
تے صرف اک ہی
طریقے نال بن
سکدی اے۔ لکھتاں
نوں \emph{اکّو
طرح نال پڑھن
جوگ} بناؤن والی
استقرائی تعریف
دا مطلب ایہ اے
کہ اسیں اوہناں
لکھتاں نوں
معنی وی اوسے
طریقے نال، یعنی
استقرائی تعریف
نال، دے سکدے آں۔

\end{document}
